\chapter{Механика частиц, четырёхимпульс и энергия покоя}

\section{Приближение, которое мы будем неоднократно использовать}

Механика релятивистской частицы должна отвечать на знакомые вопросы на новом языке. Как движется частица? Каковы её импульс и энергия? Почему эти величины объединяются в четырёхвектор, и как \(E=mc^2\) возникает из обычного аппарата механики? Прежде чем перейти к этим вопросам, нам понадобится одно элементарное приближение, к которому мы будем возвращаться на протяжении всего расчёта.

Для любого показателя степени \(p\) биномиальный ряд имеет вид
\begin{equation}
\begin{aligned}
(1+\epsilon)^p
&=1+p\epsilon\\
&\quad+\frac{p(p-1)}{2!}\epsilon^2\\
&\quad+\frac{p(p-1)(p-2)}{3!}\epsilon^3
+\cdots .
\end{aligned}
\label{eq:sr3-binomial}
\end{equation}
Показатель может быть положительным, отрицательным, дробным или даже комплексным. При достаточно малом \(\epsilon\) ряд быстро сходится, и каждый следующий член меньше предыдущего. Тогда можно остановиться на том порядке, которого требует нужная точность. Член нулевого порядка часто является наибольшим, но в приложениях ниже он нередко оказывается константой и потому динамически неинтересен. В первом поправочном члене и проявляется движение.

\begin{figure}[htbp]
\centering
\includegraphics[width=0.90\textwidth]{lecture_03_figure_01.png}
\caption{Биномиальный ряд до кубического порядка, используемый для восстановления нерелятивистской механики из релятивистских формул.}
\lectureframe{00:04:24}
\end{figure}

Обычно мы полагаем \(c=1\), поэтому скорость \(v\) означает долю \(v/c\) от скорости света. Возьмём \(\epsilon=-v^2\). При \(p=\tfrac12\)
\begin{equation}
\sqrt{1-v^2}
=
1-\frac{v^2}{2}-\frac{v^4}{8}-\cdots
\approx
1-\frac{v^2}{2},
\qquad |v|\ll 1.
\label{eq:sr3-sqrt-expansion}
\end{equation}
При \(p=-\tfrac12\) два минуса взаимно уничтожаются:
\begin{equation}
\frac{1}{\sqrt{1-v^2}}
=
1+\frac{v^2}{2}+\frac{3v^4}{8}+\cdots
\approx
1+\frac{v^2}{2}.
\label{eq:sr3-gamma-expansion}
\end{equation}
Многие множители вида \(1/2\) в классической механике задним числом можно распознать как первую поправку в одном из этих релятивистских разложений. Приближение перестаёт работать по мере того, как \(v\) приближается к 1; тогда никакое конечное усечение не воспроизводит сингулярное поведение точного выражения.

\begin{figure}[htbp]
\centering
\includegraphics[width=0.90\textwidth]{lecture_03_figure_02.png}
\caption{Обе формулы для малых скоростей видны: \(\sqrt{1-v^2}\approx1-v^2/2\) и \((1-v^2)^{-1/2}\approx1+v^2/2\).}
\lectureframe{00:07:39}
\end{figure}

\section{Что считать частицей?}

В механике частица не обязана быть элементарной. Ею могут быть электрон, Солнце, коробка шоколадных конфет или пончик. Важно лишь, чтобы объект достаточно хорошо сохранял целостность, так что его общее движение можно было описать центром масс. Пончик полезно напоминает о том, что центр масс не обязан лежать в самом материале.

В ньютоновской механике для переменных центра масс выполняются соотношения
\begin{equation}
\vec p=m\vec v,
\qquad
K=\frac12mv^2,
\qquad
\vec F=m\vec a.
\label{eq:sr3-newtonian-particle}
\end{equation}
Эти формулы не зависят от того, элементарен объект или составной. Составной объект также обладает внутренней энергией. У движущегося пончика, например, есть кинетическая энергия центра масс, а также химическая, тепловая и энергия связи:
\begin{equation}
E_{\text{Newtonian}}
=
E_{\text{internal}}+\frac12mv^2.
\label{eq:sr3-newtonian-internal}
\end{equation}
Внутренний член не зависит от общей скорости центра масс. Теория относительности сохранит это различие, но покажет, что внутренняя энергия даёт вклад в инвариантную массу всего объекта.

Необходимые инструменты — те же, что используются в классической механике: траектории, принцип стационарного действия, лагранжианы, канонические импульсы, гамильтонианы и связь между непрерывными симметриями и законами сохранения.

\section{Мировые линии и причинное разделение}

Отложим положение системы по горизонтали, а время по вертикали. Тогда её история будет кривой в пространстве-времени. Для одной частицы эта кривая — её \emph{мировая линия}. Это понятие существует и в ньютоновской механике, но в теории относительности оно становится фундаментальным, потому что преобразования Лоренца смешивают пространство и время.

Рассмотрим два события, разделённые
\begin{equation}
\Delta x^\mu=(\Delta t,\Delta\vec x),
\qquad
\Delta\vec x^{\,2}
=
(\Delta x)^2+(\Delta y)^2+(\Delta z)^2.
\end{equation}
В единицах \(c=1\) их инвариантный интервал равен
\begin{equation}
(\Delta\tau)^2
=
(\Delta t)^2-\Delta\vec x^{\,2}.
\label{eq:sr3-interval}
\end{equation}

\subsection{Времеподобное разделение}

Если \((\Delta t)^2>\Delta\vec x^{\,2}\), то \((\Delta\tau)^2>0\) и разделение времеподобно. Существует лоренцева система отсчёта, в которой оба события происходят в одной и той же пространственной точке. В этой системе соединяющая их линия вертикальна, и часы могут пройти от одного события к другому. Разность их показаний равна собственному времени \(\Delta\tau\).

\begin{figure}[htbp]
\centering
\includegraphics[width=0.90\textwidth]{lecture_03_figure_03.png}
\caption{Во времеподобном интервале временное разделение больше пространственного; существует система отсчёта, в которой оба события имеют одно и то же положение.}
\lectureframe{00:14:42}
\end{figure}

\subsection{Пространственноподобное разделение}

Если \(\Delta\vec x^{\,2}>(\Delta t)^2\), интервал пространственноподобен. Ни в одной инерциальной системе отсчёта оба события не происходят в одной и той же пространственной точке, но существует система, в которой они одновременны. Другие системы отсчёта могут обратить их временной порядок. Поэтому для пространственноподобной пары ни утверждение \emph{событие \(A\) произошло раньше}, ни утверждение \emph{событие \(B\) произошло раньше} не является инвариантным. Это и есть относительность одновременности в геометрической форме.

\begin{figure}[htbp]
\centering
\includegraphics[width=0.90\textwidth]{lecture_03_figure_04.png}
\caption{Для пространственноподобного интервала смена инерциальной системы отсчёта может сделать два события одновременными или обратить их временной порядок.}
\lectureframe{00:17:01}
\end{figure}

Массивная частица имеет времеподобную мировую линию: каждый достаточно короткий её участок удовлетворяет условию \(d\tau^2>0\). Последовательное сложение скоростей \(0.9c\) может сколь угодно приблизить скорость частицы к \(c\), но никогда не превысить его. Такое приближение напоминает парадокс Зенона. Это причинное ограничение относится не только к материальным объектам, но и к сигналам и информации, поскольку информация должна переноситься физическими системами. Сигнал быстрее света позволил бы различным системам отсчёта расходиться во мнениях о том, какой конец его истории является причиной, а какой — следствием.

\section{Четырёхскорость}

Рассмотрим короткий времеподобный участок. В обычных единицах
\begin{equation}
d\tau
=
\sqrt{dt^2-\frac{d\vec x^{\,2}}{c^2}}.
\label{eq:sr3-proper-time-dimensional}
\end{equation}
Множитель \(c^2\) диктуется размерностью: \(dt^2\) имеет размерность квадрата времени, тогда как \(d\vec x^{\,2}\) имеет размерность квадрата длины.

\begin{classroomqa}
\audiencequestion{Куда нужно вставить скорость света, если интервал записан в обычных единицах?}
\lecturerresponse{Пространственный член нужно разделить на \(c^2\), получая \(d\tau^2=dt^2-d\vec x^{\,2}/c^2\). Обычно мы вернёмся к \(c=1\), но размерностная согласованность всегда восстанавливает недостающие множители. \lecturetimestamp{00:21:23}}
\end{classroomqa}

\begin{figure}[htbp]
\centering
\includegraphics[width=0.90\textwidth]{lecture_03_figure_05.png}
\caption{Смещение вдоль мировой линии \(\Delta x^\mu=(\Delta t,\Delta\vec x)\) и соответствующее ему приращение собственного времени.}
\lectureframe{00:21:03}
\end{figure}

Обычная скорость имеет три пространственные компоненты,
\begin{equation}
v^i=\frac{dx^i}{dt},
\qquad i=1,2,3.
\end{equation}
Четырёхскорость получается делением полного смещения в пространстве-времени на инвариантное собственное время:
\begin{equation}
u^\mu
=
\frac{dx^\mu}{d\tau},
\qquad
\mu=0,1,2,3.
\label{eq:sr3-four-velocity-definition}
\end{equation}
Греческие индексы пробегают пространство-время; латинские — пространство. Нулевая компонента положения равна \(x^0=t\) в единицах \(c=1\), так что
\begin{equation}
u^0=\frac{dt}{d\tau}.
\end{equation}

\begin{figure}[htbp]
\centering
\includegraphics[width=0.90\textwidth]{lecture_03_figure_06.png}
\caption{Обычная скорость делит пространственное смещение на координатное время, тогда как четырёхскорость делит все четыре компоненты на собственное время.}
\lectureframe{00:24:05}
\end{figure}

Чтобы связать \(u^\mu\) с \(\vec v\), вынесем \(dt^2\) из выражения для интервала:
\begin{align}
d\tau^2
&=dt^2-d\vec x^{\,2}\\
&=dt^2\left(1-\frac{d\vec x^{\,2}}{dt^2}\right)
=dt^2(1-v^2).
\end{align}
Для мировой линии, направленной в будущее, берём положительный корень:
\begin{equation}
\frac{d\tau}{dt}=\sqrt{1-v^2},
\qquad
\frac{dt}{d\tau}=\frac{1}{\sqrt{1-v^2}}
\equiv\gamma.
\label{eq:sr3-time-component}
\end{equation}
Для пространственной компоненты
\begin{equation}
v^i
=
\frac{dx^i}{dt}
=
\frac{dx^i}{d\tau}\frac{d\tau}{dt}
=u^i\sqrt{1-v^2}.
\end{equation}
Следовательно,
\begin{equation}
\boxed{
u^\mu
=
\gamma(1,\vec v)
},
\qquad
u^0=\gamma,
\qquad
u^i=\gamma v^i.
\label{eq:sr3-four-velocity-components}
\end{equation}
При малой скорости \(u^0\to1\) и \(u^i\to v^i\). Вблизи скорости света и \(u^0\), и модуль пространственной части четырёхскорости становятся большими, хотя \(|\vec v|<1\).

\begin{figure}[htbp]
\centering
\includegraphics[width=0.90\textwidth]{lecture_03_figure_07.png}
\caption{Итоговая связь между обычной скоростью и всеми четырьмя компонентами \(u^\mu\).}
\lectureframe{00:30:40}
\end{figure}

\subsection{Почему независимы только три компоненты}

Четырёхскорость имеет четыре компоненты, но только три независимые степени свободы. Деление тождества для интервала на \(d\tau^2\) даёт
\begin{equation}
1
=
\left(\frac{dt}{d\tau}\right)^2
-\left(\frac{d\vec x}{d\tau}\right)^2,
\end{equation}
или
\begin{equation}
\boxed{
(u^0)^2-\vec u^{\,2}=1
}.
\label{eq:sr3-four-velocity-normalization}
\end{equation}
Одно условие нормировки устраняет одну независимую компоненту.

\begin{figure}[htbp]
\centering
\includegraphics[width=0.90\textwidth]{lecture_03_figure_08.png}
\caption{Инвариантный интервал сразу даёт условие нормировки \((u^0)^2-\vec u^{\,2}=1\).}
\lectureframe{00:33:50}
\end{figure}

\begin{classroomqa}
\audiencequestion{Можно ли понять потерю одной независимой компоненты, дифференцируя по самому собственному времени?}
\lecturerresponse{Да. Собственное время не является независимой пятой координатой; оно фиксируется пространственно-временным интервалом вдоль мировой линии. Деление этого определяющего соотношения на \(d\tau^2\) накладывает условие, чтобы четыре производные лежали на единичном времеподобном гиперболоиде \((u^0)^2-\vec u^{\,2}=1\). \lecturetimestamp{00:34:07}}
\end{classroomqa}

\section{Инвариантное действие для свободной частицы}

Какой закон определяет мировую линию при отсутствии сил? Мы сохраняем два урока классической механики:
\begin{enumerate}
\item физическая траектория делает действие стационарным, пока её концы удерживаются фиксированными;
\item действие аддитивно строится из бесконечно малых отрезков пути.
\end{enumerate}
Теория относительности добавляет третье требование: все инерциальные наблюдатели должны соглашаться относительно действия. Естественный бесконечно малый скаляр, доступный на времеподобной мировой линии, — это \(d\tau\). Поэтому действие свободной частицы имеет вид
\begin{equation}
\boxed{
S=-m\int_a^b d\tau
}
\qquad(c=1).
\label{eq:sr3-invariant-action}
\end{equation}
Запись в виде конечной суммы такова:
\begin{equation}
S=-m\lim_{N\to\infty}\sum_{k=1}^{N}\Delta\tau_k.
\end{equation}
Поскольку все наблюдатели согласны относительно каждого приращения собственного времени, все они согласны и относительно того, какая траектория делает эту сумму стационарной.

\begin{figure}[htbp]
\centering
\includegraphics[width=0.90\textwidth]{lecture_03_figure_09.png}
\caption{Действие строится как сумма инвариантных приращений собственного времени вдоль варьируемой мировой линии с фиксированными концами.}
\lectureframe{00:38:30}
\end{figure}

\begin{classroomqa}
\audiencequestion{Что происходит с энергией, если действие умножить на произвольную константу?}
\lecturerresponse{Энергия не является исходной инвариантной величиной; она будет выведена из лагранжиана и гамильтониана. Умножение всего действия на ненулевую константу не меняет его стационарных траекторий, хотя и меняет масштаб канонических величин, определяемых этим действием. Нормировка \(-m\) выбрана так, чтобы эти получающиеся величины совпадали с принятыми ньютоновскими соглашениями. \lecturetimestamp{00:39:47}}
\end{classroomqa}

\begin{classroomqa}
\audiencequestion{Почему инвариант, такой как собственное время, является хорошим кандидатом для построения действия?}
\lecturerresponse{Потому что закон стационарного действия должен выбирать одну и ту же физическую траекторию в любой инерциальной системе отсчёта. Это аналогично поиску кратчайшей кривой на евклидовой плоскости: повороты меняют компоненты координат, но не каждую бесконечно малую длину, так что они не могут изменить, какая кривая является кратчайшей. Если заменить евклидову длину собственным временем, получится релятивистская конструкция. \lecturetimestamp{00:40:39}}
\end{classroomqa}

\begin{classroomqa}
\audiencequestion{Почему для собственного времени выбирают положительный квадратный корень, а не отрицательный?}
\lecturerresponse{Собственное время определяется как положительное приращение показания часов, переносимых из прошлого в будущее. Последовательный выбор отрицательного корня лишь изменил бы знак действия, превратив минимум в максимум, но не меняя стационарную мировую линию. В системе покоя часов положительное соглашение даёт \(d\tau=dt\). \lecturetimestamp{00:42:59}}
\end{classroomqa}

Множитель \(-m\) также является соглашением о нормировке, выбранным для восстановления ньютоновской механики. После восстановления единиц обычное действие имеет вид
\begin{equation}
S=-mc^2\int_a^b d\tau.
\label{eq:sr3-action-dimensional}
\end{equation}

\begin{figure}[htbp]
\centering
\includegraphics[width=0.90\textwidth]{lecture_03_figure_10.png}
\caption{В естественных единицах конечная сумма собственного времени переходит в инвариантный интеграл \(-m\int d\tau\).}
\lectureframe{00:46:23}
\end{figure}

\section{Релятивистский лагранжиан}

Чтобы привести действие к стандартному механическому виду, вынесем \(dt\) из \(d\tau\):
\begin{equation}
d\tau
=dt\sqrt{1-\dot{\vec x}^{\,2}},
\qquad
\dot{\vec x}^{\,2}
=\dot x^2+\dot y^2+\dot z^2.
\end{equation}
Тогда
\begin{equation}
S
=
\int_a^b L\,dt,
\qquad
\boxed{
L=-m\sqrt{1-v^2}
}
\quad(c=1).
\label{eq:sr3-relativistic-lagrangian}
\end{equation}
В обычных единицах,
\begin{equation}
L=-mc^2\sqrt{1-\frac{v^2}{c^2}}.
\label{eq:sr3-relativistic-lagrangian-c}
\end{equation}
Это лагранжиан свободной частицы: он зависит от скорости, но не от положения. Ни сила, ни потенциал пока не введены.

\begin{figure}[htbp]
\centering
\includegraphics[width=0.90\textwidth]{lecture_03_figure_11.png}
\caption{После вынесения \(dt\) лагранжиан свободной релятивистской частицы выражается как функция трёх компонент скорости.}
\lectureframe{00:49:57}
\end{figure}

\subsection{Возвращение к Лагранжу и Ньютону}

Подставим \eqref{eq:sr3-sqrt-expansion}:
\begin{align}
L
&=-m\sqrt{1-v^2}\\
&\approx-m\left(1-\frac{v^2}{2}\right)\\
&=-m+\frac12mv^2.
\label{eq:sr3-lagrangian-low-speed}
\end{align}
Постоянная \(-m\) не зависит ни от положения, ни от скорости. Она вносит один и тот же вклад в любую траекторию при одних и тех же временах на концах и поэтому выпадает из уравнений Эйлера--Лагранжа. Динамически содержательная часть — это в точности
\begin{equation}
L_{\mathrm{Newton}}=\frac12mv^2.
\end{equation}
Тем самым знак и нормировка по массе в инвариантном действии выбраны правильно.

\begin{figure}[htbp]
\centering
\includegraphics[width=0.90\textwidth]{lecture_03_figure_12.png}
\caption{Разложение при малых скоростях оставляет постоянную \(-m\) и ньютоновский лагранжиан свободной частицы \(mv^2/2\).}
\lectureframe{00:53:24}
\end{figure}

Для рассматриваемой свободной частицы уравнения Эйлера--Лагранжа скажут лишь, что скорость постоянна, а мировая линия в пространстве-времени прямая. Более содержательная задача — определить сохраняющиеся импульс и энергию.

\section{Канонический импульс}

Поскольку \(L\) не зависит от положения, он инвариантен относительно пространственных переносов. Соответствующая сохраняющаяся величина — импульс. Мы не предполагаем ньютоновскую формулу \(p_i=mv_i\); мы используем каноническое определение
\begin{equation}
p_i=\frac{\partial L}{\partial\dot x^i}.
\label{eq:sr3-canonical-momentum}
\end{equation}
Для \(L=-m\sqrt{1-v^2}\),
\begin{align}
p_i
&=
-m\frac{1}{2\sqrt{1-v^2}}
\frac{\partial}{\partial\dot x^i}(1-v^2)\\
&=
-m\frac{1}{2\sqrt{1-v^2}}(-2\dot x^i)\\
&=
\frac{mv_i}{\sqrt{1-v^2}}
=m\gamma v_i.
\label{eq:sr3-relativistic-momentum}
\end{align}
Поскольку \(u^i=\gamma v^i\),
\begin{equation}
\boxed{
p^i=mu^i
}.
\label{eq:sr3-momentum-spatial}
\end{equation}
При малой скорости это сводится к \(p^i\approx mv^i\). При скоростях, близких к скорости света, знаменатель играет решающую роль.

\begin{figure}[htbp]
\centering
\includegraphics[width=0.90\textwidth]{lecture_03_figure_13.png}
\caption{Каноническое дифференцирование лагранжиана с квадратным корнем даёт \(p_x=mv_x/\sqrt{1-v^2}\).}
\lectureframe{00:58:49}
\end{figure}

\subsection{Масса, а не релятивистская масса}

Здесь \(m\) — просто масса частицы: инвариантная характеристика, которая не меняется при движении частицы. Старые тексты иногда называют её \emph{массой покоя}, а \(E/c^2\) — зависящей от скорости \emph{релятивистской массой}. Современное употребление избегает этого разделения. Движение изменяет энергию и импульс, а не инвариантную массу. Масса электрона — один и тот же параметр во всех системах отсчёта.

\section{От трёхимпульса к четырёхимпульсу}

Утверждение об одной компоненте четырёхвектора, вообще говоря, зависит от системы отсчёта. Например, \(\Delta\vec x=0\) означает, что два события имеют одно и то же положение в одной системе отсчёта, но буст наклоняет этот интервал и даёт \(\Delta\vec x'\ne0\). Напротив,
\begin{equation}
\Delta x^\mu=0
\end{equation}
инвариантно: если все четыре компоненты обращаются в нуль в одной системе отсчёта, то все четыре обращаются в нуль в любой системе отсчёта. Аналогично, \(\vec u=0\) лишь выделяет систему покоя частицы; это не инвариантное описание четырёхскорости.

Для замкнутой системы обычный закон сохранения импульса записывается как
\begin{equation}
\vec P_{\mathrm{initial}}-\vec P_{\mathrm{final}}=0.
\label{eq:sr3-three-momentum-conservation}
\end{equation}
Это уравнение для трёхвектора. Чтобы записать закон сохранения в явно лоренц-ковариантной форме, нам нужна четвёртая компонента в дополнение к трём \(p^i=mu^i\). Очевидное дополнение таково:
\begin{equation}
p^\mu=mu^\mu.
\label{eq:sr3-four-momentum}
\end{equation}
При \(c=1\),
\begin{equation}
p^\mu=(p^0,\vec p)
=
m\gamma(1,\vec v).
\end{equation}

\begin{figure}[htbp]
\centering
\includegraphics[width=0.90\textwidth]{lecture_03_figure_14.png}
\caption{Закон сохранения импульса становится инвариантным только после того, как пространственные компоненты дополняются четвёртой компонентой.}
\lectureframe{01:09:51}
\end{figure}

Недостающая компонента сама должна сохраняться. Механика частицы уже содержит ещё одну сохраняющуюся величину: энергию. Вопрос в том, является ли \(p^0=mu^0=m\gamma\) в точности этой энергией в естественных единицах. Систематический ответ даёт гамильтониан.

\section{Гамильтониан — это энергия}

Для любого не зависящего от времени лагранжиана такого вида гамильтониан равен
\begin{equation}
H=\sum_{i=1}^{3}\dot x^i p_i-L.
\label{eq:sr3-hamiltonian-definition}
\end{equation}
Эта формула не является специально ньютоновской. Это общее преобразование Лежандра, переходящее от скоростей к каноническим импульсам, и оно в равной мере применимо к релятивистскому лагранжиану.

\begin{figure}[htbp]
\centering
\includegraphics[width=0.90\textwidth]{lecture_03_figure_15.png}
\caption{Общая формула для гамильтониана: \(H=\sum_i\dot x^i p_i-L\).}
\lectureframe{01:11:55}
\end{figure}

Подставим \(p_i=m\dot x_i/\sqrt{1-v^2}\) и \(L=-m\sqrt{1-v^2}\):
\begin{align}
H
&=
\frac{mv^2}{\sqrt{1-v^2}}
+m\sqrt{1-v^2}\\
&=
\frac{mv^2+m(1-v^2)}{\sqrt{1-v^2}}\\
&=
\boxed{\frac{m}{\sqrt{1-v^2}}}
=m\gamma.
\label{eq:sr3-hamiltonian-result}
\end{align}
Именно сокращение здесь является ключевым шагом. Сравнивая с \(u^0=\gamma\),
\begin{equation}
H=mu^0=p^0.
\end{equation}
Следовательно, четвёртая компонента — это энергия:
\begin{equation}
p^\mu=(E,\vec p),
\qquad
E=H
\qquad(c=1).
\label{eq:sr3-four-momentum-natural}
\end{equation}

\begin{figure}[htbp]
\centering
\includegraphics[width=0.90\textwidth]{lecture_03_figure_16.png}
\caption{Два члена в гамильтониане приводятся к общему знаменателю; вклады \(mv^2\) сокращаются.}
\lectureframe{01:14:28}
\end{figure}

\begin{classroomqa}
\audiencequestion{Является ли формула \(H=\sum_i\dot x^i p_i-L\) только классической формулой для гамильтониана?}
\lecturerresponse{Это общее преобразование Лежандра, связанное с лагранжианом, а не специальное нерелятивистское предположение. Вид \(L\) определяет, какая механика получается. Здесь релятивистский лагранжиан с квадратным корнем даёт \(H=m/\sqrt{1-v^2}\). \lecturetimestamp{01:15:21}}
\end{classroomqa}

Инвариантная формулировка закона сохранения теперь такова:
\begin{equation}
\sum_{\mathrm{initial}}p^\mu
=
\sum_{\mathrm{final}}p^\mu,
\label{eq:sr3-four-momentum-conservation}
\end{equation}
и она объединяет сохранение энергии и всех трёх компонент импульса в одном четырёхвекторном уравнении.

\section{Энергия покоя и ньютоновский предел}

В естественных единицах точная энергия равна
\begin{equation}
E=\frac{m}{\sqrt{1-v^2}}.
\end{equation}
Используя \eqref{eq:sr3-gamma-expansion},
\begin{equation}
E
\approx
m+\frac12mv^2+\frac38mv^4+\cdots .
\label{eq:sr3-energy-natural-expansion}
\end{equation}
Второй член — привычная кинетическая энергия. Первый — энергия, которая остаётся, когда центр масс покоится.

\begin{figure}[htbp]
\centering
\includegraphics[width=0.90\textwidth]{lecture_03_figure_17.png}
\caption{Точная энергия разлагается как \(m+mv^2/2+\cdots\) при \(c=1\).}
\lectureframe{01:18:02}
\end{figure}

Восстановим размерности в точном выражении. Отношение под квадратным корнем должно быть \(v^2/c^2\), а числитель должен иметь размерность энергии:
\begin{equation}
\boxed{
E
=
\gamma mc^2
=
\frac{mc^2}{\sqrt{1-v^2/c^2}}
}.
\label{eq:sr3-energy-exact}
\end{equation}
Его разложение при малых скоростях имеет вид
\begin{equation}
\boxed{
E
=
mc^2+\frac12mv^2
+\frac38m\frac{v^4}{c^2}
+\cdots
}.
\label{eq:sr3-energy-dimensional-expansion}
\end{equation}
В покое,
\begin{equation}
\boxed{E_{\mathrm{rest}}=mc^2}.
\label{eq:sr3-emc2}
\end{equation}

\begin{figure}[htbp]
\centering
\includegraphics[width=0.90\textwidth]{lecture_03_figure_18.png}
\caption{Восстановление размерностей превращает постоянный член в \(mc^2\), после которого идёт ньютоновская кинетическая энергия \(mv^2/2\).}
\lectureframe{01:19:16}
\end{figure}

Это относится к любому замкнутому объекту, а не только к элементарной частице. \emph{В покое} означает, что центр масс покоится. Молекулы, атомы и поля внутри объекта всё ещё могут двигаться и взаимодействовать. Их внутренняя кинетическая энергия, потенциальная энергия и энергия связи вносят вклад в полную энергию покоя и тем самым в инвариантную массу объекта.

\begin{classroomqa}
\audiencequestion{Мог ли коэффициент пересчёта быть равен \(c^2\), делённому на некоторое произвольное число?}
\lecturerresponse{Одни только размерности определяют степень \(c\), но коэффициент фиксируется результатом в естественных единицах. Формула в обычных единицах должна сводиться к \(E=m\) при \(c=1\), поэтому восстановленное выражение в покое есть \(E=mc^2\) с единичным коэффициентом. \lecturetimestamp{01:22:08}}
\end{classroomqa}

\begin{classroomqa}
\audiencequestion{Имеет ли физическое значение первый член, равный единице, в биномиальном разложении, и что означают старшие члены?}
\lecturerresponse{Да. В разложении энергии первый член становится энергией покоя \(mc^2\). Следующий член — ньютоновская кинетическая энергия, а все более высокие члены — релятивистские поправки. Для бейсбольного мяча они пренебрежимо малы, но становятся существенными, когда \(v\) приближается к \(c\). \lecturetimestamp{01:22:37}}
\end{classroomqa}

Точные значения энергии и импульса расходятся при \(v\to c\):
\begin{equation}
\begin{aligned}
E&=\gamma mc^2\longrightarrow\infty,\\
\vec p&=\gamma m\vec v\longrightarrow\infty
\qquad(m>0).
\end{aligned}
\end{equation}
Следовательно, достижение скорости света потребовало бы неограниченно больших энергии и импульса. Конечная сила, действующая в течение конечного времени, сообщает лишь конечный импульс силы, поэтому массивная частица может приближаться к \(c\), но не может её достичь.

\begin{classroomqa}
\audiencequestion{Что теперь обеспечивает барьер скорости света?}
\lecturerresponse{Именно точный четырёхимпульс. И энергия, и импульс расходятся, когда массивная частица приближается к \(c\), тогда как сохранение четырёхимпульса ограничивает всякое взаимодействие. Ранее введённая нормировка четырёхскорости выражает то же утверждение в кинематической форме. \lecturetimestamp{01:24:27}}
\end{classroomqa}

\section{Масса составной системы}

Предположим, что две частицы движутся с равными по модулю и противоположно направленными импульсами. Их суммарный трёхимпульс равен нулю, так что мы уже находимся в системе покоя этой составной системы. Инвариантная масса \(M\) этой системы определяется её полной энергией в данной системе отсчёта:
\begin{equation}
Mc^2=E_{\mathrm{tot}}
\qquad
(\vec P_{\mathrm{tot}}=0).
\end{equation}
В более общем случае,
\begin{equation}
M^2c^4
=
E_{\mathrm{tot}}^2-c^2\vec P_{\mathrm{tot}}^{\,2}.
\label{eq:sr3-composite-mass}
\end{equation}

\begin{classroomqa}
\audiencequestion{Для двух частиц, движущихся в противоположных направлениях, больше ли масса составной системы, чем сумма их индивидуальных масс?}
\lecturerresponse{Если они свободны, да: в системе центра импульса их кинетическая энергия добавляется к полной энергии покоя, так что \(M>m_1+m_2\). Энергия связи при притяжении, напротив, может уменьшить полную массу, давая \(M<m_1+m_2\). Инвариантное определение — это полная энергия в системе отсчёта, где суммарный импульс равен нулю. \lecturetimestamp{01:25:06}}
\end{classroomqa}

Вот почему при взвешивании составного объекта нельзя игнорировать внутреннее движение и энергию связи. Инвариантная масса относится ко всей замкнутой системе, а не к простому арифметическому перечислению масс составляющих.

\begin{classroomqa}
\audiencequestion{Какую роль играет член энергии покоя в фазе квантового состояния?}
\lecturerresponse{Состояние с определённой энергией эволюционирует с фазой \(\exp(-iEt/\hbar)\), поэтому покоящаяся частица даёт вклад \(\exp(-imc^2t/\hbar)\). Общая для всего состояния фаза ненаблюдаема, но относительные фазы между компонентами с разными энергиями могут наблюдаться. Это квантовое следствие согласуется с приведённым здесь классическим выводом, но не требуется для него. \lecturetimestamp{01:26:57}}
\end{classroomqa}

\section{Соотношение энергия--импульс}

Нормировка четырёхскорости сразу даёт инвариантную норму четырёхимпульса. В единицах \(c=1\),
\begin{equation}
(u^0)^2-\vec u^{\,2}=1.
\end{equation}

Умножая на \(m^2\) и используя \(p^\mu=mu^\mu\),
\begin{equation}
\boxed{
E^2-\vec p^{\,2}=m^2
}
\qquad(c=1).
\label{eq:sr3-mass-shell-natural}
\end{equation}
Эквивалентно,
\begin{equation}
E=\sqrt{\vec p^{\,2}+m^2}.
\end{equation}

\begin{figure}[htbp]
\centering
\includegraphics[width=0.90\textwidth]{lecture_03_figure_19.png}
\caption{Умножение нормировки четырёхскорости на \(m^2\) даёт \(E^2-\vec p^{\,2}=m^2\).}
\lectureframe{01:31:48}
\end{figure}

Восстанавливая \(c\),
\begin{equation}
\boxed{
E^2
=
\vec p^{\,2}c^2+m^2c^4
},
\qquad
E=\sqrt{\vec p^{\,2}c^2+m^2c^4}.
\label{eq:sr3-mass-shell}
\end{equation}
Эта формула выражает энергию как функцию импульса без обращения к скорости. Она является релятивистским аналогом ньютоновского соотношения
\begin{equation}
K=\frac{\vec p^{\,2}}{2m}.
\end{equation}

\begin{figure}[htbp]
\centering
\includegraphics[width=0.90\textwidth]{lecture_03_figure_20.png}
\caption{Соотношение на массовой оболочке с восстановленными размерностями: \(E=\sqrt{\vec p^{\,2}c^2+m^2c^4}\).}
\lectureframe{01:32:38}
\end{figure}

\subsection{Безмассовые частицы}

Формула скорости \(E=\gamma mc^2\) принимает неопределённую форму \(0/0\), если \(m=0\) и \(v=c\). Поэтому скорость плохо подходит в качестве признака для различения безмассовых частиц: любая свободно распространяющаяся в вакууме безмассовая частица движется со скоростью \(c\), но разные фотоны могут иметь разные энергии.

Соотношение энергия--импульс имеет гладкий безмассовый предел:
\begin{equation}
m=0
\quad\Longrightarrow\quad
\boxed{E=c|\vec p|}.
\label{eq:sr3-massless-energy}
\end{equation}
Это относится к фотонам и, в идеализированной безмассовой теории, к гравитонам. Это превосходное приближение для ультрарелятивистской частицы всякий раз, когда \(mc^2\ll E\), в том числе для нейтрино при достаточно высокой энергии, несмотря на их малые ненулевые массы.

\begin{figure}[htbp]
\centering
\includegraphics[width=0.90\textwidth]{lecture_03_figure_21.png}
\caption{Безмассовый предел соотношения энергия--импульс даёт \(E=c|\vec p|\).}
\lectureframe{01:33:48}
\end{figure}

\section{Распад позитрония}

Рассмотрим позитроний, связанное состояние электрона и позитрона. Его инвариантная масса немного меньше \(2m_e\): составляющие имеют положительную внутреннюю кинетическую энергию, но отрицательная энергия связи больше по модулю. Численно,
\begin{equation}
m_ec^2\approx0.511\ \mathrm{MeV},
\qquad
m_e\approx9.11\times10^{-31}\ \mathrm{kg},
\end{equation}
но кинематика не зависит от численного значения.

Рассмотрим двухфотонный канал распада в системе покоя позитрония:
\begin{equation}
\mathrm{Ps}\longrightarrow\gamma+\gamma.
\end{equation}
Изначально суммарный импульс равен нулю. Поэтому закон сохранения импульса требует, чтобы фотоны вылетали в противоположных направлениях с равными по модулю и противоположными импульсами:
\begin{equation}
\vec p_1=\vec p,
\qquad
\vec p_2=-\vec p.
\label{eq:sr3-positronium-momentum}
\end{equation}
Если \(M\) — масса позитрония, то из сохранения энергии получаем
\begin{equation}
Mc^2
=E_1+E_2
=c|\vec p_1|+c|\vec p_2|
=2c|\vec p|.
\end{equation}
Следовательно,
\begin{equation}
\boxed{
|\vec p|=\frac{Mc}{2},
\qquad
E_1=E_2=\frac{Mc^2}{2}
}.
\label{eq:sr3-positronium-result}
\end{equation}

\begin{figure}[htbp]
\centering
\includegraphics[width=0.90\textwidth]{lecture_03_figure_22.png}
\caption{Покоящееся состояние позитрония распадается на два фотона с равными по модулю и противоположными импульсами; \(Mc^2=2c|\vec p|\).}
\lectureframe{01:40:55}
\end{figure}

У этого процесса нет ньютоновского аналога, в котором сохранялась бы сумма масс частиц. Начальное состояние имеет ненулевую инвариантную массу, тогда как каждый вылетающий фотон безмассов. Сохраняется полный четырёхимпульс. Запасённая энергия покоя превратилась в распространяющуюся электромагнитную энергию. Вместе с квантовым соотношением \(|\vec p|=h/\lambda\) тот же результат определяет длину волны фотона; соответствующая энергия фотона лежит намного выше оптического диапазона.

\begin{classroomqa}
\audiencequestion{Как сохраняется момент импульса при распаде?}
\lecturerresponse{Исходящее излучение должно нести тот же полный момент импульса, что и начальное состояние. В двухфотонном канале спиральности фотонов и их орбитальное состояние складываются должным образом. Точнее, спин-синглетный парапозитроний распадается преимущественно на два фотона, тогда как спин-триплетный ортопозитроний распадается преимущественно на три; разрешённый канал соблюдает все квантовые числа. \lecturetimestamp{01:42:11}}
\end{classroomqa}

\begin{classroomqa}
\audiencequestion{Почему частица, движущаяся со скоростью света, должна быть безмассовой?}
\lecturerresponse{При конечной ненулевой энергии множитель \(\gamma\) расходится при \(v\to c\). Инвариантное соотношение \(E^2-\vec p^{\,2}c^2=m^2c^4\) показывает, что частица со светоподобным четырёхимпульсом имеет \(m=0\). И наоборот, свободная безмассовая частица имеет светоподобную мировую линию и распространяется со скоростью \(c\). \lecturetimestamp{01:42:50}}
\end{classroomqa}

\begin{classroomqa}
\audiencequestion{Могло бы то же нейтральное состояние распасться на гравитоны вместо фотонов?}
\lecturerresponse{В принципе гравитационные конечные состояния допустимы, когда это разрешают квантовые числа, но гравитационная доля ветвления фантастически меньше электромагнитной. Кинематический учёт всё равно использовал бы тот же закон сохранения четырёхимпульса и \(E=c|\vec p|\) для каждого безмассового кванта. \lecturetimestamp{01:43:20}}
\end{classroomqa}

\subsection{Одна вершина распада}

Распад локален: мировая линия материнской частицы заканчивается в одном событии пространства-времени, и там же начинаются мировые линии вылетающих фотонов. По взаимодействиям в детекторе можно выполнить обратную экстраполяцию и восстановить эту общую вершину. Если бы два фотона были испущены в разных событиях, их продолженные назад светоподобные траектории не пересеклись бы в одной точке.

\begin{classroomqa}
\audiencequestion{Должны ли два фотона испускаться одновременно в любой системе отсчёта?}
\lecturerresponse{Они возникают в одном и том же событии пространства-времени, а не в двух пространственно разделённых событиях, чья одновременность могла бы зависеть от системы отсчёта. Все лоренцевы системы отсчёта согласны с тем, что одно событие — это одно событие. Квантовая механика придаёт материнской частице конечную неопределённость положения, поэтому восстановленная вершина представляет собой распределение, а не идеально известную классическую точку. \lecturetimestamp{01:44:01}}
\end{classroomqa}

\begin{classroomqa}
\audiencequestion{Почему норма четырёхскорости равна единице?}
\lecturerresponse{Собственное время было определено формулой \(d\tau^2=dt^2-d\vec x^{\,2}\) в единицах \(c=1\). Деление на \(d\tau^2\) даёт \((dt/d\tau)^2-(d\vec x/d\tau)^2=1\), а это в точности \((u^0)^2-\vec u^{\,2}=1\). Следовательно, нормировка встроена в само определение собственного времени. \lecturetimestamp{01:45:38}}
\end{classroomqa}

\begin{classroomqa}
\audiencequestion{Даёт ли \(u^0=dt/d\tau\) более глубокую интерпретацию энергии?}
\lecturerresponse{Оно даёт полезное тождество \(E=mc^2\,dt/d\tau=\gamma mc^2\): в покое отношение координатного времени к собственному равно единице, а в движении это фактор Лоренца. Помимо того, что оно кодирует эту зависимость от скорости, само по себе тождество не добавляет отдельного динамического принципа. \lecturetimestamp{01:46:36}}
\end{classroomqa}

\section{Сохранение требует замкнутой системы}

Утверждение
\begin{equation}
\sum p^\mu_{\mathrm{initial}}
=
\sum p^\mu_{\mathrm{final}}
\end{equation}
относится к изолированной системе, содержащей всё, что обменивается энергией и импульсом. В реальных лабораториях есть посторонние электрические и магнитные поля, окружающие частицы, гравитация, опоры и детекторы. Если оставить эти степени свободы вне учёта, подсистема может проявить небольшой кажущийся дисбаланс; недостающий четырёхимпульс ушёл в окружающую среду.

\begin{classroomqa}
\audiencequestion{Может ли сумма импульсов двух фотонов совсем немного отличаться от нуля?}
\lecturerresponse{Не для полной изолированной системы. В реальном эксперименте немоделируемое поле или отдача установки могут нести этот малый остаток, так что выделенная подсистема может не выглядеть замкнутой. Как только учтены все взаимодействующие части, полный четырехимпульс сохраняется. \lecturetimestamp{01:47:51}}
\end{classroomqa}

\begin{classroomqa}
\audiencequestion{Ослабляет ли принцип неопределенности закон сохранения импульса?}
\lecturerresponse{Нет. У квантового состояния положение или импульс могут быть неопределенными, и чем точнее известно одно, тем шире распределение другого, но симметрия относительно пространственных переносов все равно дает точное сохранение полного импульса для изолированной системы. Квантовая неопределенность относится к результатам и распределениям, а не к нарушению закона сохранения. \lecturetimestamp{01:49:19}}
\end{classroomqa}

\section{Точная и приближённая энергия}

Важно не путать точную формулу с её разложением при малых скоростях:
\begin{align}
E_{\mathrm{exact}}
&=
\frac{mc^2}{\sqrt{1-v^2/c^2}},
\label{eq:sr3-exact-repeated}\\
E_{\mathrm{slow}}
&=
mc^2+\frac12mv^2
+O\!\left(\frac{mv^4}{c^2}\right).
\label{eq:sr3-approx-repeated}
\end{align}
Первое выражение остается верным при любом \(v<c\). Второе полезно только при \(v/c\ll1\).

\begin{classroomqa}
\audiencequestion{Формула для энергии была лишь приближением?}
\lecturerresponse{\(E=mc^2+mv^2/2+\cdots\) — приближение. Точный результат — \(E=mc^2/\sqrt{1-v^2/c^2}\). Фактор \(v^2/c^2\) делает знаменатель безразмерным, а числитель \(mc^2\) задает единицы энергии. \lecturetimestamp{01:49:56}}
\end{classroomqa}

\subsection{Снова о картине пространства-времени}

Сделанный позже набросок на доске вновь возвращает нас к геометрии эксперимента. До распада покоящийся позитроний следует вдоль вертикальной времениподобной мировой линии. В событии распада эта мировая линия оканчивается, и из него в противоположных пространственных направлениях выходят две светоподобные мировые линии фотонов.

\begin{figure}[htbp]
\centering
\includegraphics[width=0.90\textwidth]{lecture_03_figure_23.png}
\caption{Мировая линия позитрония оканчивается в одном событии распада, из которого выходят две светоподобные мировые линии фотонов.}
\lectureframe{01:53:48}
\end{figure}

\begin{classroomqa}
\audiencequestion{Что определяет ось, вдоль которой вылетают два фотона?}
\lecturerresponse{Для неполяризованного состояния существует квантовое угловое распределение, и при отдельном распаде направление случайно выбирается из этого распределения. Как только одно направление реализовалось, сохранение импульса фиксирует фотоны строго в противоположных направлениях; саму эту ось оно заранее не выбирает. \lecturetimestamp{01:53:55}}
\end{classroomqa}

\section{Что говорят свободные уравнения движения}

Инвариантное действие было выбрано не угадыванием закона силы. Его выбрали как простейший локальный лоренц-скаляр, воспроизводящий ньютоновский лагранжиан свободной частицы при малых скоростях. Варьирование дает
\begin{equation}
\frac{d}{dt}\left(\frac{\partial L}{\partial\dot x^i}\right)=0,
\end{equation}
откуда
\begin{equation}
\frac{dp_i}{dt}=0.
\end{equation}
При фиксированной ненулевой массе это равносильно постоянству четырехскорости:
\begin{equation}
\frac{du^\mu}{d\tau}=0.
\end{equation}
Следовательно, мировая линия является прямой линией в пространстве-времени. Это одно геометрическое утверждение содержит и постоянство скорости, и неизменность пространственного направления.

\begin{classroomqa}
\audiencequestion{Существует ли более конструктивное обоснование релятивистского действия, аналогичное выбору классического действия так, чтобы оно воспроизводило уравнения Ньютона?}
\lecturerresponse{Для свободной частицы лоренц-инвариантность, локальность вдоль мировой линии и ньютоновский предел малых скоростей определяют простейшее действие с точностью до общего множителя. Его уравнение Эйлера--Лагранжа затем дает постоянство импульса, в точности ньютоновский результат при отсутствии силы. Взаимодействия дают дополнительные проверки; позже заряженная частица будет связана с электрическими и магнитными полями. \lecturetimestamp{01:54:27}}
\end{classroomqa}

\section{Два независимых классических предела}

Нерелятивистский предел и классический предел квантовой механики — разные операции. Релятивистские поправки определяются отношением \(v/c\). Квантовые эффекты определяются масштабами действия по сравнению с \(\hbar\), а также такими свойствами, как масса, размер системы и то, насколько плавно меняются внешние силы.

\begin{classroomqa}
\audiencequestion{Совпадает ли предел малых скоростей, возвращающий ньютоновскую механику, с тем пределом, в котором квантовая механика становится классической?}
\lecturerresponse{Нет. Релятивистский предел касается \(v/c\); переход от квантового к классическому касается величины \(\hbar\) по сравнению с характерным действием и часто эффективно реализуется при больших массах или плавно меняющихся силах. Система может быть квантовой, но нерелятивистской, релятивистской, но по существу классической, и тем и другим, или ни тем, ни другим. \lecturetimestamp{01:57:32}}
\end{classroomqa}

\begin{classroomqa}
\audiencequestion{Есть ли несовместимость между специальной теорией относительности и квантовой механикой?}
\lecturerresponse{Нет. Релятивистская квантовая теория непротиворечива и приводит к квантовой теории поля. Широко известная нерешённая проблема связана с включением гравитации в квантовую механику, а не с объединением специальной теории относительности с квантовыми принципами. В нынешнее обсуждение свободной частицы гравитация не входила. \lecturetimestamp{01:59:01}}
\end{classroomqa}

Пока что результат для движения без сил завершен: массивная частица следует по прямой времениподобной мировой линии, ее инвариантная масса фиксирует норму ее четырехимпульса, а энергия и импульс являются четырьмя компонентами единого сохраняющегося объекта. Введение заряда и электромагнитных полей станет следующим шагом за пределами этого свободного движения.
